\section{Conclusion}
\label{sec:conclusion}

When the unit of evaluation is the laundering attempt, the numbers AMLworld produces are
governed first by choices that papers rarely state and only then by the method. The scope rule
is worth 0.35--0.38 ring-F1. The standard split's tail inflates edge metrics and halves the
effective seed budget in the two days that carry ordinary traffic. And 44\,\% of in-scope
rings cannot be matched as connected objects by any method, because the split and the window
sever them.
% source: docs/FINDINGS.md (C9, C1 and dated notes on the window-matched budget and component-level scope)

Within those limits, a reference pipeline shows what is and is not gained. Growing rings
outward from the detector's strongest edges recovers members the detector never flags (0.5\,\%
to 25\,\% with leak-free policy inputs, at an added-member precision of 0.819; 21\,\% and 0.926 with
the day-folded inputs first reported), and ring-F1 rises from \ci{0.521}{0.429}{0.607} to \ci{0.589}{0.502}{0.675}; on
HI-Small a two-line rule matches that recall, and the learned policy adds purity. Community
detection on the same flagged edges does no better than connected components. Partitioning is
not significant at 174 rings; on HI-Medium, with 2{,}187 rings, an effect of the same size is
resolved. A 9.4-point better edge detector recovers no more rings at the same budget.
% source: results/tables/r9_ring_ci.csv; r4_partition.csv; r9_significance.csv; r6_detector_swap.csv; results/tables/v40_external_baselines.json (bar_B); results/tables/v3_closeout_p3.csv (global, full vs b2)

The practical consequence is for how group-level AML results are reported. A ring-level number
is interpretable only with its scope rule, split, window, budget allocation and matching
threshold attached, and edge F1 is not evidence about it. We release the protocol, the tables
behind every number in this paper, and the code.
